\documentclass[11pt, a4paper]{article}

% ---------- Fonts ----------
\usepackage[T1]{fontenc}
\usepackage[utf8]{inputenc}
\usepackage[dutch]{babel}
\usepackage{helvet}
\renewcommand{\familydefault}{\sfdefault}

% ---------- Page layout ----------
\usepackage[a4paper, margin=2cm]{geometry}
\setlength{\parindent}{0pt}
\setlength{\parskip}{0.55em}

% ---------- Packages ----------
\usepackage{amssymb}
\usepackage{enumitem}
\usepackage{listings}
\usepackage{xcolor}
\usepackage{booktabs}
\usepackage{array}
\usepackage{fancyhdr}
\usepackage{titlesec}
\usepackage{tcolorbox}
\usepackage{hyperref}

% ---------- Headers ----------
\pagestyle{fancy}
\fancyhf{}
\fancyhead[L]{\footnotesize\itshape Validatie-run protocol}
\fancyhead[C]{\footnotesize Datum: \rule[-0.5em]{2.5cm}{0.4pt}}
\fancyhead[R]{\footnotesize \thepage}
\renewcommand{\headrulewidth}{0.4pt}

% ---------- Section spacing ----------
\titlespacing*{\section}{0pt}{1.8ex}{0.8ex}
\titlespacing*{\subsection}{0pt}{1.2ex}{0.5ex}
\titlespacing*{\subsubsection}{0pt}{0.8ex}{0.3ex}
\titleformat{\section}{\large\bfseries}{\thesection}{0.6em}{}
\titleformat{\subsection}{\normalsize\bfseries}{\thesubsection}{0.5em}{}
\titleformat{\subsubsection}{\small\bfseries\itshape}{}{0em}{}

% ---------- Custom ----------
\newcommand{\cbox}{$\square$\;}
\newcommand{\fillline}{\rule[-0.5ex]{4cm}{0.4pt}}
\newcommand{\fillshort}{\rule[-0.5ex]{2cm}{0.4pt}}
\newcommand{\filllong}{\rule[-0.5ex]{6cm}{0.4pt}}
\newcommand{\fillblock}[1]{\rule[-0.5ex]{#1}{0.4pt}}

% ---------- Code style ----------
\lstdefinestyle{gd}{
  basicstyle=\ttfamily\footnotesize,
  backgroundcolor=\color{gray!8},
  frame=single,
  framesep=4pt,
  showstringspaces=false,
  breaklines=true,
  language=Python,
  morekeywords={var,func,const,if,else,return},
  literate={°}{{$^\circ$}}1
}
\lstset{style=gd}

% ---------- Boxes ----------
\newtcolorbox{notebox}{colback=yellow!10, colframe=yellow!50!black, boxsep=4pt, left=6pt, right=6pt, top=4pt, bottom=4pt}
\newtcolorbox{fillbox}[1][]{colback=white, colframe=black!50, boxsep=4pt, left=6pt, right=6pt, top=4pt, bottom=4pt, #1}

% ---------- Title ----------
\title{\vspace{-1.2cm}\textbf{Validatie-run protocol}\\\large stairlift live 3D telemetry}
\author{Nout Mulder \textbullet{} DeVi Comfort}
\date{}

\begin{document}
\maketitle
\thispagestyle{fancy}
\vspace{-0.7cm}

\begin{fillbox}
\textbf{Run-info} (invullen vóór je begint):

Datum: \fillline\hfill Tijd start: \fillshort\hfill Tijd eind: \fillshort

Locatie / klant: \filllong

Lift serial / ID: \fillline\hfill Software-versie: \fillshort

Tester: \fillline\hfill Begeleider aanwezig? \cbox ja \cbox nee
\end{fillbox}

\section*{Overzicht}

Eén calibratie-run, één rail, vier metingen + twee observaties. Doel:
elk DV1--DV3 criterium in de validatie-tabel van het verslag van een
echt cijfer voorzien in plaats van \texttt{[meet]}.

\begin{table}[h]
\small
\begin{tabular}{@{}lll@{}}
\toprule
\textbf{ID} & \textbf{Meting / observatie} & \textbf{Voor} \\
\midrule
T1 & Bluetooth round-trip tijd (run-gemiddelde, ms)        & DV1(b) \\
T2 & Raw encoder jitter, stationair, 10\,s (pulses)        & DV3(b) baseline \\
T3 & Kalman-filtered jitter, stationair, 10\,s             & DV3(b) \\
T4 & Positie-error bij referentiepunten (mm)               & DV2(b) \\
O1 & Visuele smoothness (geen jumps, oscillation, snapping)& observation \\
O2 & Anomalieën (BT drops, checksum errors, glitches)      & observation \\
\bottomrule
\end{tabular}
\end{table}

Pass/fail logica:

\begin{tabular}{@{}lll@{}}
\textbf{DV1(b)} pass: & T1 $\leq$ 100\,ms \\
\textbf{DV2(b)} pass: & max(T4) $\leq$ 5\,mm op rail $\leq$ 5\,m; $\leq$ 10\,mm op langere rail \\
\textbf{DV3(b)} pass: & T2 / T3 $\geq$ 10 (factor tien jitter-reductie) \\
\end{tabular}

\section{Materiaal-checklist}

\begin{itemize}[leftmargin=2em, label=$\square$]
  \item Toegang tot één werkende fysieke stairlift
  \item Meetlint (rolmaat) met mm-graduatie, minimaal 5\,m lang
  \item Schilderstape of stickertjes voor 4--5 markers
  \item Laptop met de app, Bluetooth-adapter werkend
  \item Telefoon voor foto's van het scherm bij elke marker
  \item Pen + dit document uitgeprint
  \item Reservebatterijen / oplader voor laptop
  \item Tijd: $\sim$1\,uur op de lift
\end{itemize}

\section{Code-instrumentatie (vóór de run, op kantoor)}

Drie kleine toevoegingen in de Godot-codebase. Test ze eerst in demo-mode
voordat je naar de lift gaat.

\subsection{Logfile redirect}

Start de app vanuit terminal zodat alle \texttt{print()} output naar een
bestand gaat:

\begin{lstlisting}
cd ~/MergeIPC/GD-UP-Studio
godot --path . > ~/logs/run_$(date +%Y%m%d_%H%M).txt 2>&1
\end{lstlisting}

\textbf{Bestandsnaam van jouw logfile:} \filllong

\subsection{Per-frame CSV-log (voor T2, T3)}

In \texttt{assets/code/ipc/ipc\_lift\_controller.gd}, in de functie
\texttt{on\_encoder\_received(encoder: int)}, vlak vóór de
\texttt{\_kf\_predict(dt)} aanroep:

\begin{lstlisting}
# [VALIDATIE] tijdelijk: per-meting log voor jitter-analyse
print("[CSV] %d,%d,%.6f,%.6f,%.6f" % [
    Time.get_ticks_msec(), encoder,
    _kf_pos, _kf_vel, _display_pos if not is_nan(_display_pos) else 0.0
])
\end{lstlisting}

Geeft regels: \texttt{[CSV] timestamp\_ms, raw\_encoder, kf\_pos, kf\_vel, display\_pos}.

\subsection{Bluetooth round-trip log (voor T1)}

In \texttt{assets/code/ipc/net/lift/lift\_client.gd}, in de functie
\texttt{send\_and\_receive(...)}, nét na de while-loop die op de response
wacht (rond regel 137, vóór \texttt{var result := \_pending\_result}):

\begin{lstlisting}
if _pending_result != null:
    print("[BT_RTT] msg_id=%d rtt_ms=%.1f" % [msg_id, elapsed_ms])
\end{lstlisting}

\subsection{Marker-knop (voor T4)}

Bovenaan \texttt{ipc\_lift\_controller.gd} (class-velden):
\begin{lstlisting}
var _marker_was_pressed: bool = false
\end{lstlisting}

In \texttt{\_process(delta)} aan de top:
\begin{lstlisting}
# [VALIDATIE] markeer huidige positie als referentiepunt-passage
if Input.is_key_pressed(KEY_SPACE) and not _marker_was_pressed:
    _marker_was_pressed = true
    print("[MARKER] t=%d encoder=%d kf_pos=%.5f display_pos=%.5f" % [
        Time.get_ticks_msec(), last_encoder_value,
        _kf_pos, _display_pos if not is_nan(_display_pos) else 0.0
    ])
elif not Input.is_key_pressed(KEY_SPACE):
    _marker_was_pressed = false
\end{lstlisting}

\subsection{Sanity-check vóór vertrek}

\begin{itemize}[leftmargin=2em, label=$\square$]
  \item Build de app, geen compile errors
  \item Run in demo mode: \texttt{[CSV]} regels verschijnen continu
  \item Druk spatie: \texttt{[MARKER]} regel verschijnt
  \item Demo mode genereert geen BT, dus \texttt{[BT\_RTT]} pas op locatie zichtbaar
  \item Logfile groeit, output is leesbaar
\end{itemize}

\textbf{Sanity-check voltooid op (datum/tijd):} \filllong

\newpage
\section{Op locatie — marker setup}

\subsection{Stap 0a — Lift uit en rail meten}

\begin{itemize}[leftmargin=2em, label=$\square$]
  \item Lift uitgeschakeld, sleutel uit, hoofdschakelaar uit
  \item Meet de \emph{totale} rail-lengte met meetlint (van bottom-stootblok tot top-stootblok)
  \item Tel het aantal horizontale bends visueel
  \item Schat globaal welke segment-types (straight, up-bend, down-bend, h-left, h-right)
\end{itemize}

\textbf{Totale rail-lengte:} \fillline\ \,mm

\textbf{Aantal horizontale bends:} \fillshort\hfill \textbf{Aantal verticale bends:} \fillshort

\textbf{Omschrijving rail (bv.\ ``rechte traploop met één bocht onderaan''):}

\fillblock{\textwidth}

\fillblock{\textwidth}

\subsection{Stap 0b — Marker-posities kiezen en meten}

Kies 4--5 markers, verspreid over de rail. Plak een stukje tape op de rail
bij elke gekozen positie, en meet daarna met de meetlint vanaf het
bottom-stootblok hoeveel mm de tape zit. Noteer hieronder.

\textit{Tip: zet markers exact daar waar de chair een herkenbaar
oppervlak van de rail passeert (bv.\ een bout, een rail-overgang), zodat
je tijdens de run goed kunt zien wanneer de chair de marker bereikt.}

\renewcommand{\arraystretch}{2.0}
\begin{tabular}{@{}|c|p{6.5cm}|p{3cm}|@{}}
\hline
\textbf{ID} & \textbf{Beschrijving (waar zit de marker?)} & \textbf{Fysieke positie (mm)} \\ \hline
MP1 & & \\ \hline
MP2 & & \\ \hline
MP3 & & \\ \hline
MP4 & & \\ \hline
MP5 & & \\ \hline
\end{tabular}
\renewcommand{\arraystretch}{1.0}

\section{De run zelf}

\subsection{Stap 0c — Opstart en verbinding ($\sim$10 min)}

\textit{Wat je doet:} laptop aan, terminal openen, app starten met log
redirect, en controleren of telemetrie binnenkomt.

\begin{itemize}[leftmargin=2em, label=$\square$]
  \item Laptop aan, op een vlakke ondergrond binnen handbereik
  \item Terminal openen, commando uit sectie 3.1 typen (niet uitvoeren)
  \item Lift inschakelen, sleutel aan
  \item Druk \emph{enter} op terminal — de app start, logfile begint te schrijven
  \item In de app: ga naar het IPC-scherm met de ghost chair view
  \item Verbind via Bluetooth met de lift
  \item Wacht tot in de console regels verschijnen met \texttt{[CSV]} en \texttt{[BT\_RTT]}
\end{itemize}

\textbf{Klok-tijd app gestart:} \fillshort

\textbf{Klok-tijd BT verbonden:} \fillshort

\begin{notebox}
\small \textbf{Probleem-oplossing:} Geen \texttt{[CSV]} regels? Controleer
of de code-aanpassing uit §3.2 is doorgevoerd. Geen
\texttt{[BT\_RTT]}? Controleer §3.3. App crash? Stop, herstart, noteer
de tijd en het type fout hieronder.
\end{notebox}

\textbf{Eventuele opstart-issues:}

\fillblock{\textwidth}\hfill\fillblock{\textwidth}

\subsection{Stap 1 — Stationaire baseline (10 s — voor T2 en T3)}

\textit{Wat je doet:} chair staat helemaal stil, je laat 10 seconden
data lopen. Tussendoor druk je twee keer spatie om begin/eind in de
logfile te markeren.

\begin{itemize}[leftmargin=2em, label=$\square$]
  \item Chair op MP1 (bottom-marker), motoren uit, chair geheel stil
  \item Wacht $\sim$5\,s zodat eventuele opstartfluctuaties verdwijnen
  \item \textbf{Klik spatie} in de terminal-window (lijn \texttt{[MARKER]} verschijnt). Dit is begin baseline.
  \item Noteer hieronder onmiddellijk de klok-tijd
  \item Wacht 10 seconden (gebruik telefoon-timer)
  \item \textbf{Klik spatie} opnieuw. Dit is eind baseline.
  \item Noteer eind-klok-tijd
\end{itemize}

\textbf{Baseline start (klok-tijd, hh:mm:ss):} \fillline

\textbf{Baseline eind (klok-tijd, hh:mm:ss):} \fillline

\textbf{Encoder-waarde tijdens baseline (van scherm aflezen):} \fillshort\ \,pulses

\textbf{Visuele observatie tijdens baseline (was er zichtbare beweging op het scherm?):}

\fillblock{\textwidth}\hfill\fillblock{\textwidth}

\subsection{Stap 2 — Bottom $\to$ Top run}

\textit{Wat je doet:} de lift rijdt langzaam van bottom naar top. Bij
elke marker druk je spatie + maak je foto + noteer je in de tabel
hieronder de getoonde positie.

\begin{itemize}[leftmargin=2em, label=$\square$]
  \item Joystick / startknop op de lift indrukken
  \item Lift begint te bewegen, vermoedelijk langzaam (firmware speed 20--50)
\end{itemize}

\textbf{Klok-tijd run B$\to$T gestart:} \fillshort

\textbf{Bij elke marker (MP1 t/m MP5) doe je:}

\begin{enumerate}[leftmargin=2.5em]
  \item \textbf{Kijken} naar de fysieke chair en de markers
  \item \textbf{Op het moment dat de chair fysiek over de marker rijdt:} druk \emph{spatie} op de laptop (genereert \texttt{[MARKER]} log-regel)
  \item Maak \emph{foto} van het scherm: zorg dat de getoonde positie zichtbaar is
  \item Lees de getoonde positie af (mm of \% of normalized 0--1, wat de UI ook toont) en noteer in de tabel hieronder
\end{enumerate}

\textbf{Tip:} druk \emph{nét voor} de marker; menselijke reactietijd is
$\sim$200\,ms en de chair is traag genoeg dat een tikje vroeg/laat niet
veel uitmaakt.

\renewcommand{\arraystretch}{1.9}
\begin{tabular}{@{}|c|p{2cm}|p{2cm}|p{2.5cm}|p{4.5cm}|@{}}
\hline
\textbf{Marker} & \textbf{Klok-tijd} & \textbf{Foto?} & \textbf{Display toont} & \textbf{Observatie (smooth? glitch?)} \\ \hline
MP1 (B$\to$T) & & \cbox ja \cbox nee & & \\ \hline
MP2 (B$\to$T) & & \cbox ja \cbox nee & & \\ \hline
MP3 (B$\to$T) & & \cbox ja \cbox nee & & \\ \hline
MP4 (B$\to$T) & & \cbox ja \cbox nee & & \\ \hline
MP5 (B$\to$T) & & \cbox ja \cbox nee & & \\ \hline
\end{tabular}
\renewcommand{\arraystretch}{1.0}

\begin{itemize}[leftmargin=2em, label=$\square$]
  \item Lift bereikt top station
\end{itemize}

\textbf{Klok-tijd lift bij top:} \fillshort

\textbf{Globale indruk B$\to$T (smooth motion, geen jumps?):}

\fillblock{\textwidth}\hfill\fillblock{\textwidth}

\subsection{Stap 3 — Top $\to$ Bottom return}

\textit{Idem als stap 2, maar omgekeerd. Levert tweede meting per
marker.}

\begin{itemize}[leftmargin=2em, label=$\square$]
  \item Joystick activeren in tegengestelde richting
\end{itemize}

\textbf{Klok-tijd run T$\to$B gestart:} \fillshort

\renewcommand{\arraystretch}{1.9}
\begin{tabular}{@{}|c|p{2cm}|p{2cm}|p{2.5cm}|p{4.5cm}|@{}}
\hline
\textbf{Marker} & \textbf{Klok-tijd} & \textbf{Foto?} & \textbf{Display toont} & \textbf{Observatie} \\ \hline
MP5 (T$\to$B) & & \cbox ja \cbox nee & & \\ \hline
MP4 (T$\to$B) & & \cbox ja \cbox nee & & \\ \hline
MP3 (T$\to$B) & & \cbox ja \cbox nee & & \\ \hline
MP2 (T$\to$B) & & \cbox ja \cbox nee & & \\ \hline
MP1 (T$\to$B) & & \cbox ja \cbox nee & & \\ \hline
\end{tabular}
\renewcommand{\arraystretch}{1.0}

\textbf{Klok-tijd lift weer bij bottom:} \fillshort

\subsection{Stap 4 — Afsluiten}

\begin{itemize}[leftmargin=2em, label=$\square$]
  \item Lift uit, sleutel uit
  \item App netjes stoppen (Ctrl+C in terminal, geen kill -9)
  \item Controleer logfile: open in editor, scroll door, verifieer aanwezigheid van \texttt{[CSV]}, \texttt{[MARKER]}, \texttt{[BT\_RTT]} regels
  \item Tape van rail halen
  \item Foto's naar laptop overzetten (bij voorkeur direct, anders binnen 24\,u)
\end{itemize}

\textbf{Aantal \texttt{[CSV]} regels in logfile (ongeveer):} \fillshort

\textbf{Aantal \texttt{[BT\_RTT]} regels:} \fillshort

\textbf{Aantal \texttt{[MARKER]} regels (verwacht: 10$\times$ + 2 baseline):} \fillshort

\subsection{Observaties opschrijven (binnen 10 min na de run!)}

Geheugen vervaagt snel. Schrijf nu op wat je zag, ook als het allemaal goed ging.

\textbf{O1 — Visuele smoothness:}

\fillblock{\textwidth}\hfill\fillblock{\textwidth}\hfill\fillblock{\textwidth}

\textbf{O2 — Anomalieën (BT drops, glitches, vreemd gedrag):}

\fillblock{\textwidth}\hfill\fillblock{\textwidth}\hfill\fillblock{\textwidth}

\newpage
\section{Post-run analyse (thuis / op kantoor)}

\subsection{T1 — Gemiddelde BT round-trip}

Commando:
\begin{lstlisting}
grep "\[BT_RTT\]" ~/logs/run_*.txt | \
  awk -F"rtt_ms=" '{print $2}' | \
  awk '{sum+=$1; n++} END {printf "T1 = %.1f ms (n=%d samples)\n", sum/n, n}'
\end{lstlisting}

\textbf{T1 resultaat:} \fillshort\ \,ms\hfill\textbf{n samples:} \fillshort

\textbf{DV1(b) pass?} \cbox ja ($\leq$ 100\,ms)\hspace{1em} \cbox nee

\subsection{T2 + T3 — Encoder jitter, stationaire baseline}

Extract de stationaire window uit het CSV-log. Eerst de hele CSV grep'en:

\begin{lstlisting}
grep "\[CSV\]" ~/logs/run_*.txt | sed 's/\[CSV\] //' > ~/logs/csv_data.csv
\end{lstlisting}

Open \texttt{csv\_data.csv} in een spreadsheet of in Python. Filter de
regels waarvan de timestamp tussen je baseline-start en baseline-eind
klok-tijd valt (zie stap 1; convert hh:mm:ss naar ticks\_msec). Bereken
binnen dat venster:

\begin{itemize}[leftmargin=2em]
  \item T2 = $\max(\text{raw\_encoder}) - \min(\text{raw\_encoder})$
  \item T3 = $\max(\text{display\_pos}) - \min(\text{display\_pos})$, omgerekend naar pulses
\end{itemize}

\textbf{T2 raw jitter (peak-to-peak):} \fillshort\ \,pulses (dus $\pm$\fillshort\ \,pulses)

\textbf{T3 Kalman jitter (peak-to-peak):} \fillshort\ \,pulse-equivalent

\textbf{Ratio T2/T3:} \fillshort\ \,$\times$

\textbf{DV3(b) pass?} \cbox ja (ratio $\geq$ 10)\hspace{1em} \cbox nee

\subsection{T4 — Positie-error per marker}

Voor elke marker: vergelijk fysieke positie (sectie 4) met getoonde
positie. Marker-passage log-regels:

\begin{lstlisting}
grep "\[MARKER\]" ~/logs/run_*.txt
\end{lstlisting}

De \texttt{display\_pos} in de log is normalized (0--1). Omrekenen naar
mm: \texttt{display\_pos $\times$ rail\_length\_mm}.

\renewcommand{\arraystretch}{2.0}
\begin{tabular}{@{}|c|c|c|c|c|c|@{}}
\hline
\textbf{Marker} & \textbf{Fysiek (mm)} & \textbf{Display B$\to$T (mm)} & \textbf{Display T$\to$B (mm)} & \textbf{Avg (mm)} & \textbf{Error $|$ mm$|$} \\ \hline
MP1 & & & & & \\ \hline
MP2 & & & & & \\ \hline
MP3 & & & & & \\ \hline
MP4 & & & & & \\ \hline
MP5 & & & & & \\ \hline
\end{tabular}
\renewcommand{\arraystretch}{1.0}

\textbf{T4 max error:} \fillshort\ \,mm\hfill\textbf{T4 gemiddelde error:} \fillshort\ \,mm

\textbf{Rail-lengte was:} \fillshort\ \,mm\hfill threshold: \cbox 5\,mm ($\leq$ 5\,m) \cbox 10\,mm (langer)

\textbf{DV2(b) pass?} \cbox ja (max $\leq$ threshold)\hspace{1em} \cbox nee

\section{Verslag invullen}

Open \texttt{project/main.tex}, ga naar \texttt{Table tab:validation}
(rond regel 1100). Vervang \texttt{[meet]} per rij:

\renewcommand{\arraystretch}{1.5}
\begin{tabular}{@{}|p{5.8cm}|p{2.5cm}|p{2.5cm}|p{3.5cm}|@{}}
\hline
\textbf{Verslag-rij} & \textbf{Bron} & \textbf{In te vullen} & \textbf{Pass/fail tekst} \\ \hline
Rail length & sectie 4 & \fillshort\ mm & --- \\ \hline
Number of horizontal bends & sectie 4 & \fillshort & --- \\ \hline
Position error at end station & MP1 of MP5 & \fillshort\ mm & DV2(b): \cbox pass \cbox fail \\ \hline
Max position error at ref. points & T4 max & \fillshort\ mm & DV2(b): \cbox pass \cbox fail \\ \hline
Encoder jitter (raw, stationary) & T2 & $\pm$\fillshort\ pulses & --- \\ \hline
Encoder jitter (after Kalman) & T3 & \fillshort\ pulses & DV3(b): \cbox pass \cbox fail \\ \hline
Bluetooth round-trip time & T1 & \fillshort\ ms & DV1(b): \cbox pass \cbox fail \\ \hline
Visual smoothness through bends & O1 & korte zin & --- \\ \hline
\end{tabular}
\renewcommand{\arraystretch}{1.0}

\textbf{Logged events paragraaf} (regel $\sim$1126): vul werkelijke
observaties O2 in.

Build PDF:
\begin{lstlisting}
cd project && latexmk -pdf main.tex
\end{lstlisting}

Commit:
\begin{lstlisting}
git add project/main.tex project/main.aux project/main.fdb_latexmk \
        project/main.fls project/main.log project/main.pdf
git commit -m "validate: fill measurement results from physical run"
\end{lstlisting}

\section{Eerlijkheids-regels}

\begin{itemize}[leftmargin=2em, label=$\square$]
  \item Bij failure (criterium niet gehaald): noteer de werkelijke waarde,
        markeer \emph{fail}, en bespreek de oorzaak in de Conclusion van
        het verslag. Niet ``een tweede run doen tot het past''.
  \item Eén run is genoeg voor proof-of-concept; N=1 is duidelijk in de
        Limitations vermeld.
  \item Bewaar de raw logfile in \texttt{logs/} (niet committen, te groot).
        Bewaar wel een geanonimiseerde samenvatting als \texttt{logs/run\_summary.txt}.
  \item Niets verzinnen. Als T1 niet gemeten is wegens ontbrekende
        instrumentatie, rapporteer dat: ``BT round-trip not measured
        in this run; instrumentation required for future validation.''
\end{itemize}

\vspace{1.5em}

\begin{fillbox}
\textbf{Afsluitende notities / lessons learned:}

\fillblock{\textwidth}

\fillblock{\textwidth}

\fillblock{\textwidth}

\fillblock{\textwidth}
\end{fillbox}

\vspace{0.5em}
\footnotesize Versie van dit protocol: \today. Plak een ingevulde kopie
bij het verslag op de map / sla op in \texttt{logs/}.

\end{document}
